\sect{一、单项选择题（每小题 2 分，共 20 分）}
\begin{enumerate}[label=\arabic*.]
\item \judge 某特征在训练集的取值为 $2,4,6,8$。用训练集的均值和总体标准差做 z-score，用训练集的最小值和最大值做 min-max 变换。两种变换都不截断，取值为 $10$ 的测试样本分别变为多少？
\opts{$\sqrt 5,\ 1$}{$\sqrt 5,\ 4/3$}{$\sqrt {15}/2,\ 4/3$}{$\sqrt 2,\ 1$}

\item \judge 同一测试集上，模型甲的预测绝对误差为 $0,0,6$，模型乙的预测绝对误差为 $3,3,3$。分别按 MAE 和 MSE 选择误差较小的模型，两次各选哪个？
\opts{甲、甲}{乙、甲}{乙、乙}{甲、乙}

\item \judge 某树的一个子树在独立验证集上错分 $3$ 个样本，用该节点训练样本的多数类代替整棵子树后错分 $2$ 个样本。后剪枝规则为验证错分数下降时执行。哪项同时正确说明三种经典建树准则和本题的处理？
\optss{ID3 用信息增益，C4.5 在信息增益筛选后比较增益比，分类 CART 可比较子节点的加权基尼指数，本题应剪枝}{ID3 用增益比，C4.5 用信息增益，分类 CART 用加权基尼指数，本题应剪枝}{ID3、C4.5 和分类 CART 都直接用训练错分率选属性，本题应剪枝}{三种建树准则的描述与 A 项相同，但因训练子树的信息量变少，本题应保留子树}

\item \judge 特征 $X$ 为离散变量。取 $X=x$ 时，与类别 1、2、3 的联合概率依次为 $0.03$、$0.09$、$0.08$。判错的代价为 $1$，判对为 $0$。求最优预测类别及给定 $X=x$ 的最小错误概率。
\opts{类别 $2$，$0.45$}{类别 $3$，$0.55$}{类别 $2$，$0.55$}{类别 $2$，$0.11$}

\item \judge 考虑满足经典大样本条件的 1-NN 分类器，记其极限风险为 $R_\infty$，贝叶斯风险为 $R^\ast$。另考虑均匀分布于单位超立方体的数据。下列说法正确的是
\optss{对任意有限训练集都保证测试风险不超过 $2R^\ast$，维数越高越容易得到足够多的近邻}{大样本风险必等于贝叶斯风险，且固定半径的超立方邻域质量与维数无关}{大样本风险满足 $R^\ast\le R_\infty\le 2R^\ast$。在超立方体中心，半径固定为 $0<r<1/2$ 时，$\ell_\infty$ 球概率质量为 $(2r)^d$，会随维数 $d$ 增大而下降}{大样本风险下界为 $2R^\ast$，固定相同半径时高维邻域总覆盖更多概率质量}

\item \judge 某坐标早期出现过一次很大的梯度，之后很多步的梯度都较小。下列对 AdaGrad 和 RMSProp 的解释正确的是
\optss{AdaGrad 记录梯度的一阶和，利用正负号抵消早期大梯度的影响}{AdaGrad 累计平方梯度，早期大值长期留在累计量中。RMSProp 对平方梯度做指数移动平均，较早贡献会随时间衰减}{两种方法都只使用当前平方梯度，历史梯度不会影响分母}{RMSProp 对平方梯度逐步等权累加，所以早期大梯度会一直以相同权重被计入}

\item \judge 三份卷积模型笔记分别写着“以堆叠的小卷积核构建深网络”“在一层中并联多个尺度的分支，再拼接通道”“学习残差并加上恒等捷径（维度匹配时）”。它们最符合哪组经典架构，次序与笔记一致？
\opts{VGG、GoogLeNet、ResNet}{GoogLeNet、VGG、ResNet}{ResNet、GoogLeNet、VGG}{VGG、ResNet、GoogLeNet}

\item \judge 某 KNN 分类器的学习曲线趋于平稳，训练与验证错误率都接近 $0.28$。同一任务上，经独立检验已有模型达到约 $0.08$ 的错误率。当前 KNN 的 $k$ 较大。哪项解释或操作最有依据？
\optss{训练与验证误差接近，说明模型已达到任务能取得的最佳结果}{学习曲线主要表明方差过高，应继续增大 $k$，无须再检验}{只需不断扩充训练样本，既有误差必然降到 $0.08$}{训练误差也高且与验证误差接近，符合偏差偏大的情形。较大的 $k$ 可能过度平滑，可在验证中检验适当减小 $k$ 是否有利}

\item \judge 一个三分类器对每条记录预测一个类别。一个验证集中，三类各有 $80,15,5$ 个真实样本，其中分别判对 $72,6,1$ 个。先计算每类召回率再取等权平均，称为宏平均。先合计各类分子、分母再取比值，称为微平均。宏平均和微平均召回率依次是多少？
\opts{$0.79,\ 0.79$}{$0.50,\ 0.79$}{$0.79,\ 0.50$}{$0.50,\ 0.50$}

\item \judge 二值信道的输入 $X$ 等概率取 $0$ 或 $1$。输出 $Y$ 也取 $0$ 或 $1$，且 $P(Y=1\mid X=0)=0$，$P(Y=1\mid X=1)=1/2$。采用以 $2$ 为底的对数。$H(Y\mid X)$、$H(X\mid Y)$、$I(X;Y)$ 依次约为多少？
\opts{$0.6887,\ 0.5000,\ 0.3113$}{$0.8113,\ 1.0000,\ 0.3113$}{$0.5000,\ 0.6887,\ 0.1887$}{$0.5000,\ 0.6887,\ 0.3113$}
\end{enumerate}

\sect{二、判断对错题（每小题 1 分，共 10 分）}
\begin{enumerate}[label=\arabic*.,start=11]
\item \judge 对只有两个连续特征的决策树，若每个内部节点只对一维特征作阈值测试，则决策区域内部边界的各段都在与坐标轴平行的直线上。
\item \judge 为评估模型在新患者上的表现，只要训练集和测试集没有完全相同的影像文件，同一患者在不同时间采集的影像分到两个集合也不会引入信息泄露。
\item \judge 在只保留第一主成分的 PCA 中，把第一单位特征向量 $\bm u$ 改成 $-\bm u$ 后，中心化样本投影得分改变，所以用同一新主方向反变换得到的重构也会改变。
\item \judge 单独看一个步长为 $1$、无填充的 $1\times1$ 卷积层，每个输出位置只读取输入同一空间位置，但可用多个输入通道的信息计算一个输出通道。
\item \judge 高斯核 Parzen 窗估计中，每个样本核都经过带宽归一化，在实数全集上的积分为 $1$。扩大带宽会改变密度曲线的形状，但估计密度在实数全集上的积分仍为 $1$。
\item \judge 在 HMM 中，由于隐状态的一阶马尔可夫性质，$P(i_t\mid o_1,\ldots,o_t,\lambda)$ 必与 $P(i_t\mid o_1,\ldots,o_T,\lambda)$ 相同，其中 $t<T$。
\item \judge 固定数据和 $\mathrm{MinPts}$，增大 DBSCAN 的 $\varepsilon$ 可能使原来的核心点失去核心点身份。
\item \judge 若将 Lasso 的平方误差和改为原来的 $1/n$，为了让新目标函数与旧目标函数仅相差一个正常数倍，原来的 $L_1$ 惩罚系数也应除以 $n$。
\item \judge 对无位置编码、无掩码的全局自注意力，重新排列输入 token 的顺序，输出会按相同的顺序重新排列。
\item \judge 用同一组参数计算一个含 $2$ 个样本的小批量和一个含 $6$ 个样本的小批量的平均梯度，将两组平均梯度各乘以 $1/2$ 相加，总能得到这 $8$ 个样本的平均梯度。
\end{enumerate}

\sect{三、计算与分析题（共 8 题，共 70 分）}

\noindent 如无另行规定，$\log$ 表示自然对数，计算结果可保留精确式或写到 4 位有效数字。

\noindent\textbf{1. 小批量分类网络的反向传播与归一化}\pts{10}

用两个样本做三分类训练。输入 $\bm{x}_A=(1,0)^{\mathsf T}$、$\bm{x}_B=(0,1)^{\mathsf T}$，分别属于第 $2$ 类和第 $1$ 类。网络对任意输入计算
\[
h=\sigma(\bm{w}^{\mathsf T}\bm{x}+b),\qquad
\bm{z}=\bm{v}h+\bm{c},\qquad \bm{p}=\operatorname{softmax}(\bm{z}).
\]
取 $r=\ln 2$，$\bm{w}=(\ln 3,-\ln 3)^{\mathsf T}$、$b=0$、$\bm{v}=(4r,0,-4r)^{\mathsf T}$、$\bm{c}=(-2r,0,2r)^{\mathsf T}$。损失 $J$ 是两个样本的交叉熵的算术平均，对数以 $e$ 为底。写出有理数或含 $r$ 的精确结果即可。

\begin{enumerate}[label=(\arabic*)]
\item 求两个样本的隐藏层输出。\pts{1}
\item 求各样本的三个类别概率，并求小批量损失 $J$。\pts{1}
\item 对每个样本，写出其单样本损失对三个输出得分的梯度。\pts{1}
\item 利用第（3）问的结果，求 $J$ 对 $\bm{v},\bm{c}$ 的梯度。\pts{1}
\item 分别求两个样本的单样本损失对隐藏层激活前输入 $\bm{w}^{\mathsf T}\bm{x}+b$ 的偏导。\pts{1}
\item 求 $J$ 对 $\bm{w},b$ 的梯度。\pts{1}
\item 用本批次的梯度做一次动量更新。对每个参数使用 $u_{\rm new}=\tfrac12 u_{\rm old}+g$、$\theta_{\rm new}=\theta-\tfrac1{10}u_{\rm new}$。仅 $b$ 的旧动量非零，为 $r/4$。给出更新后的全部参数。\pts{1}
\item 另有某层的两个样本激活值组成 $\begin{pmatrix}1&3\\3&7\end{pmatrix}$，行是样本，列是特征。对该矩阵做 BatchNorm，给出均值、方差和归一化结果。方差按元素个数平均，缩放为 $1$、平移为 $0$，稳定项取 $0$。\pts{1}
\item 对第（8）问同一个原始矩阵做 LayerNorm，其他设置相同，给出均值、方差和归一化结果，并指出两种方法的统计轴有何不同。\pts{1}
\item 若改为按 $A,B$ 顺序分别进行一次单样本 SGD 更新，说明为什么这两步通常不能用开始时的平均梯度直接代替。\pts{1}
\end{enumerate}

\noindent\textbf{2. 支持向量机、异或与多项式核}\pts{8}

两个开关 $r_1,r_2\in\{0,1\}$ 恰有一个开启时，控制器输出正类，否则输出负类。令 $x_j=2r_j-1$，训练点依次为 $(-1,-1)^\top$、$(-1,1)^\top$、$(1,-1)^\top$、$(1,1)^\top$，对应标签 $(-1,+1,+1,-1)$。核为 $K(\bm x,\bm z)=(1+\bm x^\top\bm z)^2$。

\begin{enumerate}[label=(\arabic*)]
\item 利用两类样本连线的中点说明单个线性判别面无法严格分开这两类。\pts{0.5}
\item 用两个隐藏阈值单元和一个输出阈值单元实现异或。约定 $H(s)=1$ 当 $s\ge0$，否则为 0。输入用原始开关 $r_1,r_2$，输出 1 代表正类。\pts{1}
\item 展开核函数，给出一个显式映射 $\bm\varphi(\bm x)$，使 $K(\bm x,\bm z)=\bm\varphi(\bm x)^\top\bm\varphi(\bm z)$。\pts{1}
\item 按给定点序写出 $4\times4$ 的 Gram 矩阵 $G$。\pts{1}
\item 写出硬间隔对偶的核形式。利用 $G=8I+\bm 1\bm 1^\top$ 及偏置的对偶等式约束，化简目标函数。\pts{1}
\item 求四个最优对偶乘子。说明为何得到的是最优解。\pts{1}
\item 用一个支持向量求偏置 $b$，再把决策函数 $f(\bm z)$ 化成关于 $z_1,z_2$ 的多项式。\pts{1}
\item 求新点 $\bm z=(2,-1/2)^\top$ 的判别分数和类别。\pts{0.5}
\item 若给对偶乘子加上软间隔上界 $\alpha_i\le C$，求使上述硬间隔对偶解仍为软间隔最优解的最小 $C$，说明理由。\pts{1}
\end{enumerate}

\noindent\textbf{3. 混合模型的更新与聚类评价}\pts{12}

对四个二维样本 $(-1,-1),(-1,1),(1,-1),(1,1)$ 拟合两成分 GMM。第 1、2 成分的初始权重为 $4/5,1/5$，初始均值为 $(-1,0),(1,0)$，初始共享协方差为 $\sigma^2I_2$，其中 $\sigma^2=2/\ln 2$。M 步同时更新权重、均值与共享的球形方差。以下第 (1)、(2)、(6)–(11) 问使用各问所给的独立数据。
\begin{enumerate}[label=(\arabic*)]
\item 在同一个一维训练集上，用相同的 $k$ 做 KNN 密度估计。两个查询点到第 $k$ 个近邻的半径分别为 1 和 2，窗口未被边界截断，且均无距离平局。求前一个估计值与后一个估计值之比。若换用固定带宽的高斯 Parzen 窗，写出高斯核带宽过小时密度估计的典型表现。\pts{1}
\item 某 Mini-Batch K-means 的中心 $(-2,0)$ 是前四个分配样本的均值，新批次分给它的点为 $(1,1),(1,-1)$。若用累计平均，求更新后的中心。另在一轮 K-means++ 选点中，只剩 $u,v$ 两个候选点。$u$ 到当前最近中心的距离为 2，$v$ 被选中的概率为 $9/13$。反求 $v$ 到当前最近中心的距离。\pts{1}
\item 对题首 GMM，写出第 2 成分与第 1 成分的责任度之比，算出四个样本对第 2 成分的责任度。\pts{1}
\item 继续完成这一轮 EM，求两成分更新后的权重和均值。\pts{1}
\item 用第 (3)、(4) 问的结果估计更新后的共享方差 $\sigma_{\mathrm{new}}^2$。给出 M 步公式和精确值。\pts{1}
\item 另一项聚类实验的真实类别与预测簇列联表为
\[
\begin{array}{c|rrr}
 &C_1&C_2&C_3\\\hline
Y_1&3&1&0\\
Y_2&1&1&2
\end{array}
\]
求纯度和兰德指数。兰德指数按无序样本对计算。\pts{1}
\item 继续使用第 (6) 问的列联表。以“两点同簇”为预测正例，以“两点同属真实类别”为实际正例，求配对精确率、召回率和 $F_1$。\pts{1}
\item 用以 2 为底的对数计算第 (6) 问列联表的互信息 $I(Y;C)$。约定计数为零的项贡献为零。\pts{1}
\item 根据第 (8) 问的结果，求同质性、完整性和两者等权的 V-measure。继续使用以 2 为底的对数。\pts{1}
\item 某轮层次聚类中，簇 $A$ 有一个对象，簇 $B$ 有三个对象。$A$ 和 $B$ 已经合并。$A$ 中的对象到外部单点簇 $C$ 的距离是 3，$B$ 中三个对象到 $C$ 的距离为 6、8、10。求合并后的簇到 $C$ 的单链、全链和平均链距离。\pts{1}
\item 方法甲用高斯核对邻域样本加权，每步将当前位置移到加权均值。方法乙设定一个二维原型网格，每读入一个样本，就更新离它最近的原型及其在网格上的邻居。分别识别两种方法，并说明它们最终如何确定簇或组织表示。\pts{2}
\end{enumerate}

\noindent\textbf{4. 线性自编码器、隐空间表示与 VAE}\pts{8}

二维测量向量 $\bm{x}$ 已中心化，协方差矩阵为 $\bm{\Sigma}$。单位特征向量 $\bm{u}_1=(1,1)^{\mathsf T}/\sqrt2$、$\bm{u}_2=(1,-1)^{\mathsf T}/\sqrt2$ 分别对应特征值 $4$ 和 $1$。考虑一维瓶颈的线性自编码器，编码为 $z=\bm{u}^{\mathsf T}\bm{x}$，解码为 $\widehat{\bm{x}}=\bm{u}z$，且 $\|\bm{u}\|_2=1$。损失取 $\E\|\bm{x}-\widehat{\bm{x}}\|_2^2$。

\begin{enumerate}[label=(\arabic*)]
\item 写出编码和重构的维数。\pts{0.5}
\item 展开重构误差，证明期望损失等于 $\operatorname{tr}(\bm{\Sigma})-\bm{u}^{\mathsf T}\bm{\Sigma}\bm{u}$。\pts{1}
\item 求最优的编码方向和最小损失，并说明它与第一主成分的关系。\pts{1}
\item 若训练去噪自编码器，给定一对同源测量值 $(\bm{x},\widetilde{\bm{x}})$，其中 $\widetilde{\bm{x}}$ 由 $\bm{x}$ 加噪得到，训练时分别用哪个作为输入和重构目标？\pts{0.5}
\item 稀疏自编码器限制哪一层的哪些量？堆栈自编码器怎样利用前一编码层的输出？\pts{0.5}
\item 设另一个自编码器将两种状态分别编码为 $-2$ 和 $6$。从前一个编码向后一个编码线性插值，取进度 $1/4$，求插值点。若解码器为 $g(z)=z^2$，该点的解码值是否等于这两个端点解码值的线性插值？用数值说明。\pts{1}
\item 两个 VAE 在同一验证集上的平均对数重构似然和平均先验 KL 分别为 $(-3,1)$ 和 $(-5,0)$。目标为对数重构似然减去 $\beta$ 倍的先验 KL。对 $\beta>0$，比较两者何时得分更高、何时相同。\pts{1}
\item 结合第（7）问的第二个模型，说明平均先验 KL 为零对编码器意味着什么。为判断解码器是否忽略潜变量，还可做什么检查？\pts{1}
\item 对一个新样本，取 $x=3$，解码分布为 $x\mid z\sim\mathcal N(z^2,1)$。令 $z=\mu+s\epsilon$，$s>0$。固定此次标准正态噪声样本 $\epsilon=1/2$，在 $\mu=1/2,s=1$ 处，求单样本对数重构似然 $\ln p(x\mid z)$ 对 $\mu,s$ 的偏导。\pts{1.5}
\end{enumerate}

\noindent\textbf{5. 空洞卷积、门控网络与 Transformer 解码器}\pts{8}

\begin{enumerate}[label=(\arabic*)]
\item 输入特征图高 $7$、宽 $9$。一层 $3\times3$ 卷积在两方向的空洞率均为 $2$，步长均为 $1$，不填充。求等效核的空间跨度和输出的高、宽。\pts{0.5}
\item 对两个通道的特征图 $\begin{pmatrix}2&0\\4&2\end{pmatrix}$ 和 $\begin{pmatrix}1&3\\-1&9\end{pmatrix}$ 做全局平均池化，写出输出向量。\pts{0.5}
\item 双向 RNN 在某位置给出前向状态 $(1,2)$ 和后向状态 $(-1,3)$。若按前向在前、后向在后的顺序拼接，写出该位置的表示。解释用这个表示做实时预测时，为何可能用到尚未到达的输入。\pts{1}
\item 某标量 GRU 使用 $h_t=(1-z_t)h_{t-1}+z_t\widetilde h_t$，$\widetilde h_t=\tanh(x_t+r_t h_{t-1})$。已知 $h_{t-1}=1$，重置门、更新门的激活前得分分别为 $\ln 3,0$，两门均使用 Sigmoid。取 $x_t=\tfrac12\ln 3-\tfrac34$。依次求 $r_t,z_t,\widetilde h_t,h_t$。\pts{1.5}
\item 对 LSTM，设某一步的输入门、遗忘门、候选记忆均固定。若把该步的输出门置零，说明该步的隐状态及细胞状态各怎么变。\pts{0.5}
\item 在 Seq2Seq 翻译任务中，目标输出是依次为 $u,v,\mathrm{EOS}$ 的三个词元。训练时采用 teacher forcing，首个解码输入为 $\mathrm{BOS}$。写出三次解码输入。推理时若第一步预测为 $w\ne u$，第二次解码输入是什么？\pts{1}
\item 一个带注意力的 Seq2Seq 解码器在本步为三个编码状态 $(2,0),(0,2),(2,2)$ 分配权重 $1/2,1/3,1/6$。求本步上下文向量，并说明下一步的上下文是否必须与本步相同。\pts{0.5}
\item Transformer 的源序列有 $5$ 个位置，当前批量计算 $3$ 个目标位置，每位置有一条模型表示。做单头交叉注意力时，查询、键维数为 $2$，值维数为 $3$。按“位置为行”的约定，说明 $Q,K,V$ 分别来自哪一侧，并写出 $Q,K,V$、得分矩阵及加权输出的尺寸。\pts{1}
\item Transformer 解码层按掩码自注意力、交叉注意力、前馈网络的先后顺序计算。设进入本层的目标序列表示为 $X$，编码器的输出为 $H$，每个子层均在残差相加后做 LayerNorm。依次写出三个子层的更新式，可以用 $\operatorname{MaskedAttn}$、$\operatorname{CrossAttn}$ 和 $\operatorname{FFN}$ 表示三个子层。\pts{1}
\item 在第（9）问的前馈网络子层中，改变一个位置的输入，能否直接改变其他位置在该子层的输出？说明原因。\pts{0.5}
\end{enumerate}

\noindent\textbf{6. 实验记录分类与判别面}\pts{8}

实验记录有 $A,B$ 两类，先验分别为 $2/3,1/3$。词表依次是“标注、模拟、校准”。$A$ 有四篇文本，三个词的累计出现次数为 $(5,3,2)$，包含三个词的文档数为 $(4,2,2)$。$B$ 有两篇文本，相应的两组计数为 $(1,1,4)$ 和 $(1,1,2)$。文本特征均采用加一平滑。
\begin{enumerate}[label=(\arabic*)]
\item 文本仅包含两次“模拟”。用多项式朴素贝叶斯模型求后验赔率 $P(A\mid x)/P(B\mid x)$，保留分数。\pts{1}
\item 如果删去其中一次“模拟”，第 (1) 问的后验赔率会变成多少？比较删词前后的赔率，说明重复词如何改变多项式模型的证据。\pts{1}
\item 改用伯努利朴素贝叶斯，求“模拟、模拟”在两类下的条件概率，再算后验赔率。说明删去一次“模拟”是否影响这个赔率。\pts{1}
\item 另用一个类别特征“来源”单独建模，来源的三个可能水平为甲、乙、丙。$A$、$B$ 两类计数分别为 $(3,1,0)$、$(1,0,1)$。新记录来源为丙。分别算出未平滑和加一平滑的 $P(A\mid\text{来源为丙})$。\pts{1}
\item 现改用两个连续特征。$A$ 类有 4 条样本，两坐标的样本和为 $(0,8)$、平方和为 $(4,32)$。$B$ 类有 2 条，相应两组统计量为 $(4,-2)$ 和 $(10,10)$。用极大似然法求两类的均值和方差向量。\pts{1}
\item 给定类别时，假设两个连续特征条件独立并分别服从高斯分布。利用第 (5) 问估计值，化简 $\ln(P(A\mid x)/P(B\mid x))$，写出两类的等后验判别面方程。\pts{1}
\item 根据第 (6) 问的模型，求坐标为 $(1,0)$ 的新记录属于 $A$ 的后验概率，并给出分类。\pts{1}
\item 只将 $B$ 类第二个特征的方差改为 8，其他量均保持原值。求此时后验对数优势中 $x_2^2$ 的系数，并说明一般情况下等后验判别面是否仍为直线。\pts{1}
\end{enumerate}

\noindent\textbf{7. 多项式回归的稀疏估计与多类逻辑回归}\pts{8}

为估计设备的标定曲线，采用 $\hat y=w_0+w_1x+w_2x^2$，仅对 $w_1,w_2$ 加正则。假设训练目标独立服从均值为 $\hat y$、方差为 $\sigma^2$ 的高斯分布。给两个斜率系数独立的拉普拉斯先验 $p(w_j)=(2s)^{-1}\exp(-|w_j|/s)$，$s>0$，截距不加先验惩罚。本题的 $\log$ 均为自然对数。

\begin{enumerate}[label=(\arabic*)]
\item 写出该模型含截距的特征向量。哪一个系数取零时，该模型关于 $x$ 至多为一次？\pts{0.5}
\item 分别写出 MLE 和 MAP 关于参数的优化目标，用似然、先验表达即可，并说明两者在使用先验上的区别。\pts{1}
\item 将 MAP 目标写成 $\frac12\sum_i(y_i-\hat y_i)^2+\lambda(|w_1|+|w_2|)$ 的最小化问题，求 $\lambda$ 与 $\sigma^2,s$ 的关系。\pts{1}
\item 在另一组已处理的数据上，固定截距后平方损失恰为 $\frac12[(w_1-2)^2+(w_2-0.5)^2]$。取 $\lambda=1$，分别检查零点和非零半轴，求正则化后的 $(w_1,w_2)$。\pts{1}
\item 画出或描述 $|w_1|+|w_2|\le1$ 的形状，从约束边界解释 Lasso 为何容易得到零系数。\pts{1}
\item 另一任务用三类 softmax 回归。某个样本 $\bm x=(2,-1)^\top$ 的 logit 向量为 $\bm a=(\ln 3,\ln 2,0)^\top$，真实标签是第 2 类。求预测概率向量与单样本交叉熵。\pts{1}
\item 从 $L=\ln\sum_j e^{a_j}-a_2$ 出发，求通用的 $\partial L/\partial a_k$ 的表达式。\pts{1}
\item 若 $\bm a=W\bm x+\bm b$，$W$ 的第 $k$ 行对应类别 $k$，求（6）这个样本对 $W$ 第二行的梯度。\pts{0.5}
\item 给该样本的三个 logit 同时加常数 $c$，概率是否改变？用 softmax 公式说明。\pts{0.5}
\item 写出 softmax 回归中 $\ln(p_1/p_2)$ 与 $\bm x$、参数行向量和偏置的关系，并在（6）的点上求值。\pts{0.5}
\end{enumerate}

\noindent\textbf{8. 浮标观测的缩放前向计算}\pts{8}

某浮标所在海域有三种隐藏状态，记为 $1,2,3$，每次信号读数为 $X$ 或 $Y$。用 HMM $\lambda=(A,B,\pi)$ 建模。$A$ 的行、列分别是当前和下一状态，$B$ 的列按 $X,Y$ 排列。
\[
 A=\frac14\begin{pmatrix}1&2&1\\1&1&2\\2&1&1\end{pmatrix},\qquad
 B=\frac14\begin{pmatrix}2&2\\1&3\\3&1\end{pmatrix},\qquad
 \pi=\tfrac14(1,2,1),\qquad O=(X,Y,X,X).
\]
每一步从前一时刻已归一化的前向量出发，先算待归一化向量 $\widetilde\alpha_t$，其分量之和记为 $c_t$，归一化后的向量记为 $\widehat\alpha_t$。以下可保留分数，数值对数用自然对数。
\begin{enumerate}[label=(\arabic*)]
\item 写出 $t=1$ 和 $t>1$ 时 $\widetilde\alpha_t$ 的计算式，以及用 $c_t$ 得到 $\widehat\alpha_t$ 的公式。\pts{1}
\item 求 $c_1$ 和 $\widehat\alpha_1$。\pts{1}
\item 求 $c_2$ 和 $\widehat\alpha_2$。\pts{1}
\item 求 $c_3$ 和 $\widehat\alpha_3$。\pts{1}
\item 求 $c_4$ 和 $\widehat\alpha_4$。\pts{1}
\item 用四个缩放因子表示并计算 $P(O\mid\lambda)$ 与 $\ln P(O\mid\lambda)$。\pts{0.5}
\item 定义缩放后向量 $\widehat\beta_4=(1,1,1)$，并按
\[
 \widehat\beta_t(i)=\frac{\sum_j a_{ij}b_j(o_{t+1})\widehat\beta_{t+1}(j)}{c_{t+1}}
\]
递推。求 $\widehat\beta_3$，进而求 $P(S_3=3\mid O,\lambda)$。\pts{1}
\item 给定已经收到的四个读数，用 $\widehat\alpha_4$ 计算下一个读数为 $X$ 的条件概率。\pts{0.5}
\item 说明为什么长序列计算时要逐步缩放并累计对数。就隐藏状态的表示及更新机制而言，说明标准离散 HMM 与普通 RNN 的一项关键区别。\pts{1}
\end{enumerate}

\vfill \begin{center}\zihao{5} —— 试题结束 ——\end{center}
